\chapter{State of the Art}
\label{cha:state-of-the-art}

Research on operational diagnosis addresses several different tasks. Some
systems reconstruct the requests passing through an application. Others
rank suspicious metrics, explain security events or use agents to investigate
and repair a failure. These tasks are related, but success at one does not
establish success at the others. This chapter compares their evidence and
uses that comparison to motivate design and evaluation choices for the
Norlys proof of concept.

The chapter begins with a focused comparison of conventional ML and
LLM-based diagnosis. The systematic review then contributes 68 retained
reports from IEEE Xplore, ACM Digital Library, Google Scholar and two
separately recorded citation additions. Appendix~\ref{app:systematic-literature-review}
explains the search, selection and appraisal. The synthesis prioritises the
strength of diagnostic evidence and the evaluation choices it supports;
it does not rank systems by incomparable headline scores.

The tracing, conventional ML and security studies also provide technical
background for interpreting the operational path. Such contextual sources
are not automatically members of the LLM-focused review corpus. Publication
versions and study-specific limitations are recorded in the appendix and
assessment workbook.

\section{Conventional ML and LLM-Based Diagnosis}
\label{sec:state-of-the-art:ml-llm}

LLMs are themselves machine-learning (ML) models. The distinction here is
between methods trained or fitted for a specific diagnostic task and
workflows that use a pretrained LLM to interpret evidence and guide an
investigation. This focused comparison examines representative approaches
to establish the rationale for the LLM-focused review. It considers their
inputs, diagnostic outputs and evaluation limits, rather than providing
an exhaustive review of conventional ML.

\subsection{Conventional Methods for Operational Diagnosis}

Conventional ML can do more than detect unusual behaviour. Sage learns
relationships between service dependencies, resource measurements and
latency without requiring manually labelled root causes. It uses
counterfactual analysis: estimating whether performance would recover if
a suspected service or resource behaved normally. Its experiments report
root-cause identification accuracy above 93\% for the evaluated performance
problems. However, the method still needs representative observations of
the system. The authors discuss using fault injection when ordinary
operation provides too few failure modes
\parencite[Sections~3--4, 6--7]{gan2021sage}.
A shortage of labelled incidents therefore does not, by itself, rule out
an ML approach.

CIRCA offers another option. It fits regression models to measurements
before a fault is detected and uses a structural graph to assess which
metric changes are possible root-cause indicators. This reduces its
dependence on examples of previously diagnosed failures. Its requirements
instead include suitable pre-failure measurements and knowledge of the
system's architecture and metric relationships
\parencite[Section~4]{li2022circa}. These assumptions matter when deciding
whether a method can be transferred to a different environment.

Conventional models can also combine telemetry types. Eadro learns from
logs, metrics and traces together, using service dependencies to support
anomaly detection and rank suspected services. It trains against anomaly
and culprit-service labels. Its evaluation uses injected CPU and network
faults in TrainTicket and SocialNetwork, with similar root-cause
distributions in the training and test sets
\parencite[Sections~IV--V]{lee2023eadro}. This demonstrates that using
several telemetry sources is not a capability unique to LLMs. However,
ranking a service does not itself provide an explanation of which
configuration caused the fault or why a particular repair is appropriate.

\subsection{LLM-Based Investigation and Its Limits}

RCAgent illustrates a different use of a model. A controller selects tools
to retrieve information and consult specialised log and code analysis
agents. It can request additional observations as an investigation
develops. The surrounding software provides system access, retains
observations and checks tool calls
\parencite[Sections~2--3]{wang2024rcagent}. The attraction for this thesis
is the possibility of combining investigation steps with an explanation
that refers to the retrieved evidence. This capability depends on the
tools and information supplied to the model. It does not give the model
automatic knowledge of the deployed system.

RCACopilot provides a direct comparison with conventional classifiers.
On 653 Microsoft incidents, using a 75/25 training--test split, its GPT-4
workflow reports a micro-F1 of 0.766, compared with 0.076 for FastText and
0.022 for XGBoost. Its simpler GPT-4 prompting variant scores 0.026.
These results support the value of the complete workflow in that setting,
including diagnostic information and historical examples. They do not
isolate model type as the reason for the difference or establish superiority
over conventional ML generally
\parencite[Sections~5.1--5.3, Table~2]{chen2024rcacopilot}.

The boundary between the approaches is also flexible. RCACopilot uses
FastText embeddings trained on historical incidents to retrieve relevant
examples for the LLM \parencite[Section~4.2]{chen2024rcacopilot}. A system
can therefore combine specialised learning methods with language-model
reasoning. For the PoC, a similar division could let statistical or learned
methods identify suspicious observations while an agent investigates their
meaning and requests further evidence.

LLM-based investigation introduces its own failure modes. RCAgent includes
mechanisms for malformed outputs, invalid actions and repeated tool calls,
and retains original observations to manage long tool responses
\parencite[Sections~2--3]{wang2024rcagent}. These mechanisms show that a
usable agent requires more than a prompt. A generated explanation still
needs to be checked against the fault and the source records. Neither a
well-written answer nor agreement between agents establishes correctness.

\subsection{Implications for the PoC and Review Scope}

Table~\ref{tab:sota:ml-llm} summarises the comparison. The entries describe
the examined approaches and the questions they raise for this project.
They are not universal properties of every ML model or LLM system.

\begin{table}[htbp]
\centering
\small
\caption{Comparison of the examined diagnostic approaches.}
\label{tab:sota:ml-llm}
\begin{tabularx}{\textwidth}{@{}p{0.19\textwidth}XX@{}}
\toprule
Aspect & Conventional ML approaches & LLM-based workflows \\
\midrule
Data requirements &
Sage and CIRCA use system observations without labelled root causes.
Eadro uses anomaly and service labels. &
RCAgent needs tools and system context. RCACopilot also uses historical
incidents. \\
Diagnostic output &
Suspected services, resources or metrics, depending on the method. &
Cause categories or written diagnoses, potentially with evidence and
proposed actions. \\
Investigation &
Analysis follows the implemented model and diagnostic procedure. &
A tool-using controller can choose further retrieval steps from returned
observations. \\
Main checks &
Representative data, dependency assumptions and performance on faults
outside the training distribution. &
Correct interpretation, valid tool use, supported claims and consistency
across runs. \\
\bottomrule
\end{tabularx}
\end{table}

For the Norlys PoC, the intended outcome is a diagnosis of a failed request
across operational layers, supported by records that a person can inspect.
An LLM-based workflow is worth investigating because it offers a mechanism
for interpreting those records, choosing further queries and presenting
an explanation. The comparison also leaves a role for conventional ML
where a task has suitable measurements and a clearly defined output.
The amount and representativeness of available data must be examined for
each candidate method. Limited traffic alone cannot establish which method
will work best.

A meaningful comparison must keep the diagnostic task consistent. An anomaly
detector should be assessed on detection, while methods that identify causes
can be compared on the correct component and failure explanation. The
reported scores above concern different datasets and outcomes, so they
cannot be combined into a ranking. For this project, a proposed evaluation
would use the same controlled faults and comparable access to incident
evidence. It would report training or setup requirements alongside diagnosis
time, computational cost, repeated-run variation and unsupported claims.
Any training or retrieval examples must be separated from the test incidents.

This comparison motivates investigating LLM-based agents without assuming
that they are the best solution. The following literature review focuses
on how such agents use operational telemetry and how studies establish the
correctness and reliability of their diagnoses. Its findings will inform
the choice of agent approach and evaluation criteria. Conventional ML
remains relevant as a possible component or comparison method, while the
choice between a single agent and multiple agents requires separate
evidence.

\section{Telemetry Correlation and Root-Cause Analysis}
\label{sec:state-of-the-art:correlation}

Reconstructing a request means identifying which observations describe
its execution. Diagnosing its failure adds a further question: which
component or condition explains the observed outcome? TraceWeaver addresses
the first task by matching incoming and outgoing spans using timing,
call-graph constraints and statistical optimisation. Its evaluation uses
microservice applications and Alibaba traces. It also demonstrates how
approximate reconstructed traces can help locate delayed services
\parencite[Sections~4--6]{ashok2024traceweaver}.

TraceWeaver is relevant when shared identifiers are unavailable, but it
requires more than a collection of loosely related log messages. Its
implementation obtains request and response information and learns
constraints from test traces. Its production-data evaluation separates call
graphs and compresses inter-arrival times to increase concurrency. The
results consequently support a particular reconstruction method under
specified observations, rather than arbitrary reconstruction across
unobserved network layers. For the PoC, an explicit identifier match should
be distinguishable from a link inferred from time, endpoint and topology.
The latter needs an uncertainty statement and supporting observations.

Sage and CIRCA investigate causes using system dependencies. Sage combines
an RPC graph with performance observations and counterfactual analysis.
Its experiments include resource interference and software faults in
microservice applications. Its incomplete-instrumentation experiment still
retains per-tier latency and end-to-end measurements. It should therefore
not be read as evidence that diagnosis remains reliable when those
observations disappear entirely
\parencite[Sections~5--6]{gan2021sage}.

CIRCA ranks root-cause metrics using a structural graph and regression-based
hypothesis testing. In its evaluation on 99 Oracle database incidents,
AC@1 is 0.404 and AC@5 is 0.763. Here AC@k measures recovery of labelled
root-cause metrics among the first k recommendations. These values are not
an accuracy estimate for identifying an Azure resource or a firewall rule.
Constructing CIRCA's structural graph also requires knowledge of the system
architecture and the meaning of its metrics
\parencite[Sections~5.1--5.4, Table~3]{li2022circa}.

An independent empirical comparison provides a useful qualification.
Pham, Ha and Zhang evaluate causal discovery and root-cause methods on
synthetic data and microservice testbeds. They find difficulties in
recovering edge directions, increasing computational cost with larger
systems and sensitivity to the assumed failure time. In their implementation,
CIRCA's original structural graph is replaced by a graph learned using the
PC algorithm because the required call graph is unavailable. The comparison
therefore also shows how changing an input assumption changes the method
being evaluated \parencite[Sections~3--5]{pham2024howfar}.

For RQ1, these studies motivate an explicit record of the relationships
that can actually be observed. For RQ3, they motivate testing diagnosis with
missing records, uncertain incident times and competing symptoms. They do
not establish that a learned causal graph is necessary for the PoC. A
simpler evidence-based baseline is needed before attributing an improvement
to a more complex method.

\section{Security Telemetry and Feature Selection}
\label{sec:state-of-the-art:security}

Security observations need context. A connection count has a different
meaning when it describes blocked connections to an internal service than
when it describes routine outgoing traffic. SIERRA makes this distinction
explicit by grouping SIEM events by severity, network direction and action
context. It derives counts, rates, host and port diversity, and distribution
features, then ranks anomalous time windows for investigation. Its evaluation
uses approximately 29.4 million events from two firewalls and an intrusion
detection system \parencite[Sections~4--7]{lee2022sierra}.

SIERRA offers a concrete basis for considering derived features in RQ2,
but its aggregation settings should not be copied without evaluation. It
uses ten-minute windows and rolling training periods of eight weeks. Its
labels are assembled from analyst investigations, threat intelligence and
additional checks, rather than exhaustive labelling of every event. A
low-traffic PoC may provide neither a comparable history nor enough events
for the same features to be informative. In that setting, the rule applied
to one failed request may be more useful than a population-level anomaly
score.

Feature quality also depends on how the underlying data was produced.
Engelen, Rimmer and Joosen reconstruct and examine CICIDS2017. They identify
problems in flow construction, attack labelling and features that allow
shortcut learning. Their corrected datasets change performance for specific
attack classes while aggregate accuracy remains approximately 0.99. This is
evidence for checking ground truth and reporting results by scenario, not
for treating a high aggregate score as sufficient validation
\parencite[Sections~III--IV]{engelen2021troubleshooting}.

The Holmes preprint connects security telemetry to agent investigation.
It uses counters and sampled traffic for triage, then constructs compact
packet-evidence summaries for an LLM. Reports must refer to observations in
those summaries. However, the representative incident table contains
incorrect attributions with high reported confidence. Traceable evidence
makes such errors inspectable without guaranteeing the conclusion is
correct. Its replay-based evaluation also does not establish production
performance or the correctness of automated mitigation
\parencite[Sections~3--4, Table~2]{chen2026holmes}.

Together, these findings suggest two separate checks for a proposed
security feature: can it be computed reliably from the available telemetry,
and does it help distinguish the controlled scenarios? Candidate fields
for the Norlys path include time, source and destination, request target,
status, policy action and rule identifier. This is a proposed requirements
list for the PoC, not a claim that every field has been validated by the
review or is already available from every source. Retaining the source
records allows an investigator to check an aggregate or model-generated
interpretation against the underlying event.

\section{LLM Agents for Operational Diagnosis}
\label{sec:state-of-the-art:agents}

Operational LLM systems differ in how they obtain information. RCACopilot
uses incident handlers to gather diagnostic information and retrieve
historical examples before predicting a root-cause category. Its evaluation
uses 653 incidents from one Microsoft service, split into training and test
sets. The reported GPT-4 result is a micro-F1 of 0.766 and a macro-F1 of
0.533. The distinction matters because the categories are imbalanced. Its
ablation also finds that adding more incident-related information does not
necessarily improve the prediction
\parencite[Sections~5.1--5.4]{chen2024rcacopilot}.

RCAgent instead gives a controller access to tools and specialised LLM
experts. Its evaluation uses 161 selected anomalous Flink jobs, with
annotations reviewed by engineers. It assesses generated causes, solutions
and evidence, as well as tool-use behaviour. An ablation using more general
query tools produces substantially more invalid actions. This supports
careful tool design, but does not isolate agent count as the sole reason
for improvement. The system changes tool interfaces, output validation,
observation handling and agent organisation together
\parencite[Sections~4--5]{wang2024rcagent}.

The two evaluations also measure different outcomes. RCACopilot's category
F1 cannot be ranked directly against RCAgent's text-similarity measures or
LLM-assigned scores. A plausible explanation may resemble a reference
answer while pointing to the wrong component. Conversely, a correct answer
may use different wording. For the PoC, correctness should be assessed
against the injected fault and separately from the relevance and accuracy
of the cited evidence.

The TimeRAG preprint illustrates another measurement distinction. It
combines time-series representations with retrieval and an LLM. Its
multiple-choice diagnosis accuracy and its LLM-judged report win rate
produce different pictures of performance. On SockShop, for example,
TimeRAG's reported accuracy is lower than the strongest baseline in the
same table, although its report comparisons are favourable. The work also
notes limitations from time-series augmentation and offline evaluation
\parencite[Sections~4--6, Tables~2--3]{park2026timerag}.

Live benchmarks expose additional difficulties. AIOpsLab distinguishes
fault detection, localisation, root-cause analysis and mitigation. Its
experiments show that performance varies across these tasks, and that
additional steps can increase cost without consistently improving success.
The paper describes repeated tool calls, invalid commands and difficulty
using large telemetry outputs
\parencite[Sections~3.2--3.6]{chen2025aiopslab}.

The SREGym preprint further separates diagnostic and mitigation outcomes
across 90 scenarios, with three runs per agent--model configuration. Added
background disturbances reduce diagnosis success for every evaluated
configuration. It uses a structured diagnosis rubric and reports checks
against a human evaluator. Its results support testing competing symptoms
and repeated runs. They do not establish a universal ranking of agent
architectures, since the agents also differ in models, tools and budgets
\parencite[Sections~2.5--3]{clark2026sregym}.

These studies motivate a comparison in which candidate agent approaches
receive the same accessible evidence, incident description and
time or tool budget. The particular baseline remains an evaluation-design
decision. Useful measurements include correct faulty component, correct
failure explanation, unsupported claims, invalid tool calls, elapsed time
and variation across repeated attempts. Agent agreement should not be used
as a substitute for checking the fault and its supporting evidence.


\subsection{What Agent Comparisons Actually Establish}

mABC combines specialised roles, a workflow and weighted voting. Its
TrainTicket ablation reports a root-cause score of 54.4 for the complete
GPT-4-Turbo system and 38.4 without multi-agent collaboration. This is
useful evidence for investigating collaboration. However, the score is
$\mathrm{RA}=(A_c-0.1A_i)/A_t$, where $A_c$ and $A_i$ are correct and
incorrect proposed causes and $A_t$ is the total reference causes. It is
not ordinary incident-level accuracy. The ablation also does not match
inference cost between the variants. The result therefore supports a
specific design hypothesis, not the conclusion that multiple agents are
always preferable \parencite[Section~3.2, Table~7]{zhang2024mabc}.

SuperDRI shows how a graph built from past incidents can guide investigation
of cloud databases. A coordination agent, action planner and Kusto agent
retrieve context and execute diagnostic queries. Its evaluation reports
735 successful first attempts out of 957 KQL queries and an average of
2.39 actions matching engineers' actions over 61 held-out incidents.
These measure tool execution and process agreement, respectively. Neither
establishes whether the final root cause is correct. The paper does not
define a diagnostic ground-truth scoring procedure or a matched baseline
\parencite[Section~2.2, Figure~2, Table~1]{zhu2026superdri}.
For the PoC, successful Log Analytics queries are therefore an intermediate
capability. A separate check must establish whether the evidence supports
the diagnosis.

DBAgent combines telemetry analysis, retrieval and a trained reasoning
policy. Its preprocessing includes statistical anomaly detection and
rule-based execution-plan analysis, illustrating a hybrid use of specialised
methods and an LLM. Its success criterion jointly checks a root-cause label
and a remedy; the overall results also include a consultation task. These
results must not be interpreted as root-cause-only accuracy. Moreover,
DBAgent receives supervised and reinforcement learning, whereas its main
comparison agents use prompting. The difference cannot be attributed
solely to agent organisation \parencite[Sections~4--6]{chen2026dbagent}.

Agent-MIRUPD, a preprint describing a staged single-agent design, links diagnosis stages and tests the consequences
of incorrect information passed between them. This is relevant to a
multi-agent pipeline: an early mistake can become a later agent's assumed
fact. Its comparison with published RE1 and RE2 results is explicitly
literature-based; those baselines were not rerun in the same environment.
Consequently, it is not independent confirmation of their original scores
or a controlled demonstration of superiority
\parencite[Sections~6.3--6.4]{rasouli2026agentmirupd}.

Together, these examples suggest separating three experimental changes:
the information supplied to an agent, the organisation of its reasoning
and the resources available to it. Adding agents while also changing the
model, tools and budget prevents a clear explanation of any improvement.
A single investigator and a collaborating variant can instead share a
model, evidence interface and scenario set. Their respective strengths can
then be assessed under both a common resource limit and their observed
operating cost. This is a proposed experimental design derived from the
comparison, not an experimental result of this thesis.

\section{Automated Remediation and Human Oversight}
\label{sec:state-of-the-art:remediation}

A remediation recommendation and an executed repair require different
forms of validation. RCAgent evaluates proposed solutions, whereas
AIOpsLab provides live tasks in which an agent changes the system. The
latter makes it possible to observe recovery, but benchmark success still
depends on what the evaluator checks
\parencite{wang2024rcagent,chen2025aiopslab}.

STRATUS provides the most direct assessed example of controlled
multi-agent remediation. It separates detection, diagnosis, mitigation
and undo roles. Its transactional no-regression argument relies on
serialised writers, faithful undo and checks against a severity measure.
The implementation records reversible changes and supports undo-and-retry.
The paper also acknowledges that perfect undo is difficult when changes
affect application-specific state or external systems
\parencite[Sections~3--4]{chen2025stratus}.

In the reported GPT-4o experiments, STRATUS resolves 9 of 13 AIOpsLab
mitigation problems. Removing retry reduces success to 2 of 13, while
retrying without undo succeeds on 3 of 13. This is evidence for evaluating
rollback as part of the repair process. The paper nevertheless reports
that eight ITBench successes exploit faults that stop persisting after a
pod restart. Those outcomes cannot demonstrate repair of an equivalent
persistent production fault
\parencite[Section~5, Tables~2--3]{chen2025stratus}.

STRATUS describes support for intercepting actions for human approval and
using dry-run commands where applicable. Its reported results do not
establish the effectiveness of a human approval workflow. Nor does the
transactional argument imply that intermediate states are harmless to
users. Its guarantee concerns the defined observable state sequence under
the stated assumptions. Human decision quality, approval delay and
side-effects outside the measured state require their own evaluation.

The PoC should therefore make the proposed change concrete before approval.
The operator needs its target, parameters, supporting evidence, expected
effect and rollback plan. A separate execution boundary can enforce the
permitted actions. Recovery verification should check both the selected
endpoint and the relevant configuration, rather than treating command
completion as repair. These are design implications for RQ4. Their
implementation and effectiveness remain matters for the thesis experiments.


\section{Evidence for Evaluating Diagnostic Agents}
\label{sec:state-of-the-art:evaluation}

\subsection{Define the Diagnosis Before Scoring It}

A diagnostic answer can identify a component, a fault type, an incident
time or a propagation path. These are distinct outcomes. OpenRCA makes
this explicit by requiring all elements requested in a query to match:
a partly correct answer is distinguished from a fully correct answer.
Its best reported agent result with Claude 3.5 Sonnet is 11.34\% full-query
accuracy over 335 cases. This apparently low score cannot be compared
directly with service-localisation accuracy, where identifying the faulty
component alone may suffice
\parencite[Sections~2.3--5, Tables~1--4]{xu2025openrca}.

For this thesis, the primary diagnostic outcome should require both the
faulty component or operational layer and a correct explanation of the
injected condition. Component-only results can be reported separately.
Evidence grounding is a further outcome: cited records must exist and
support the asserted relationship. A correct answer obtained from an
unsupported guess should not receive the same grounding assessment as an
answer backed by the observed request, rule or configuration.

The reference answer must come from the controlled fault and its verified
system effect, rather than from another generated diagnosis. The injected
configuration, expected symptoms and recovery checks can be recorded
before the agent runs. Scenario labels that reveal the answer should be
kept out of agent-visible evidence. This matters in the existing PoC,
where the controlled-failure log explicitly names its scenario. That
endpoint demonstrates telemetry delivery and correlation; by itself, it
is not a demanding test of root-cause reasoning.

\subsection{Test False Diagnoses and Incomplete Evidence}

Souza and Cola\c{c}o J\'unior evaluate an MCP-based agent across two fault
scenarios and one healthy control. Their best configuration reports
71.1\% success across 45 runs pooled over all three scenarios. This is
not a fault-only accuracy estimate. In the healthy control, the same
configuration produces two false diagnoses in 15 runs, or 13.3\%; the
reported 95\% interval is 3.7--37.9\%. Across the eight configurations,
false-positive rates range from 0\% to 100\%
\parencite[Tables~4--5]{souza2026cloudmicroservices}.
The small denominators also show why a percentage without its case count
and uncertainty can be misleading.

The PoC should therefore contain healthy requests as well as injected
failures. It should distinguish a false diagnosis from an explicit statement
that evidence is insufficient. Missing logs, competing symptoms and an
incorrect initial suspicion are relevant perturbations. SREGym's
background-disturbance experiments and Agent-MIRUPD's altered upstream
context support testing these conditions
\parencite{clark2026sregym,rasouli2026agentmirupd}.
An agent must not be rewarded for confidently naming a cause merely because
every evaluation case is guaranteed to contain a visible fault.

\subsection{Control Comparisons and Report the Cost}

The same incidents, accessible observations and scoring rules are needed
for an interpretable comparison. Model version, prompts, retrieval context,
tools and stopping conditions should be recorded. When comparing one and
several agents, a common token, tool-call or time limit helps distinguish
organisational effects from simply purchasing more inference. Training or
historical examples must be separated from evaluation incidents. DBAgent's
trained-versus-prompted comparison illustrates why training access must
also be reported \parencite{chen2026dbagent}.

Repeated runs assess the stability of a configuration on the same cases.
They do not turn one fault into many independent fault types. Results
should therefore report both the number of scenarios and repetitions,
with per-scenario outcomes alongside aggregate results. Timeouts, invalid
outputs and failed tool calls belong in the run accounting. Excluding them
from the denominator can make an unreliable system appear accurate.
SREGym and AIOpsLab provide examples of evaluating repeated attempts and
operational failures \parencite{clark2026sregym,chen2025aiopslab}.

A shorter response is not necessarily a cheaper investigation. Useful
resource measurements include elapsed diagnosis time, model calls, input
and output tokens, and tool calls. Model prices and deployment choices
belong with cost calculations so they can be reproduced. These are
secondary outcomes alongside correctness, not substitutes for it.

\begin{table}[htbp]
\centering
\small
\caption{Evaluation criteria derived from the review. Exact thresholds and
scenario sampling belong in the evaluation design.}
\label{tab:sota:evaluation}
\begin{tabularx}{\textwidth}{@{}p{0.24\textwidth}X@{}}
\toprule
Criterion & Operational interpretation for the PoC \\
\midrule
Diagnostic correctness & Component and cause match the injected fault; report component-only and joint results separately. \\
Evidence grounding & Cited observations exist, are relevant and support the explanation; record unsupported claims. \\
Healthy-system behaviour & Count false diagnoses and appropriate abstention on healthy or insufficient-evidence cases. \\
Robustness & Repeat cases and perturb evidence availability, background symptoms and initial assumptions. \\
Fair comparison & Share cases, evidence, model and budgets when isolating agent organisation; disclose training access. \\
Efficiency & Record time, tokens, model calls and tool calls, including failed attempts. \\
Operational reliability & Report timeouts, non-submission, invalid outputs and tool failures with all-run denominators. \\
\bottomrule
\end{tabularx}
\end{table}

These criteria provide the basis for selecting an approach through the
PoC experiments. They do not yield an architecture ranking from the
literature alone. A collaborating approach is preferable only if its
measured benefits justify its additional cost and complexity under the
chosen scenarios and constraints.

\section{Comparison and Research Gap}
\label{sec:state-of-the-art:gap}

The assessed evidence supports combining several ideas rather than
selecting one paper as a complete solution. Request reconstruction supplies
relationships between observations. Security-feature analysis supplies
context for interpreting events. Agent studies address information gathering
and diagnostic explanations. Remediation studies address state changes and
recovery. Table~\ref{tab:sota:implications} summarises the implications.

\begin{table}[htbp]
\centering
\small
\caption{Evidence-informed priorities for the PoC. These are evaluation
priorities, not completed implementation results.}
\label{tab:sota:implications}
\begin{tabularx}{\textwidth}{@{}p{0.11\textwidth}XX@{}}
\toprule
Question & Design priority & Evidence to collect \\
\midrule
RQ1 & Preserve identifiers and distinguish inferred links. & Correct and missing path links, with source records. \\
RQ2 & Select interpretable fields and test useful aggregates. & Availability, scenario relevance and effects of removing features. \\
RQ3 & Use bounded tools and check explanations against the fault. & Component and cause correctness, evidence quality, cost and repeated-run variation. \\
RQ4 & Separate approval, execution, recovery and undo. & Approved scope, actual changes, recovery tests, rejected actions and rollback outcomes. \\
\bottomrule
\end{tabularx}
\end{table}

The contribution is a controlled investigation within the selected
operational path. Combining Cloudflare and Azure alone does not establish
novelty. The assessed studies do not show that application tracing is
sufficient to explain failures in a separate security or network layer,
or that agent agreement and a successful restart establish a correct
diagnosis and safe repair.

The thesis can test these questions by making the path, injected faults,
observations and permitted actions explicit. Its claims should remain
bounded by the evaluated scenarios. The retrieval limits documented in Appendix~\ref{app:systematic-literature-review}
mean that this is a bounded empirical focus, rather than a claim that no
other study addresses these questions.
